\section{One window reduction for growing odd coefficient primes}
\label{sec:odd-prime-uniform-actual-window}

\paragraph{Exact retained scope.}
This conditional-interface argument retains every omission of
AC-G174-window and adds two disjoint families: the high two-power
characters of AC-G175 at coefficient primes at least five other
than $269$, and the pure two-power characters of AC-G177 at
coefficient three.  The numerical remainder and mixed realization
are still required.  No stronger source package from AC-G172 is
introduced, and coefficient two is not covered by the new theorems.

Keep the actual fields, page, prescribed pair and coefficients:
\[
 \begin{aligned}
 F_u&=\mathbb Q(\zeta_{2^{h(u)+2}},\zeta_7,
                                  \sqrt{-q},\sqrt{-d},\omega),\\
 H_u&=F_u(\sqrt[3]{269}),\qquad
 G_u=\operatorname{Gal}(H_u/K),\qquad K=\mathbb Q(\omega),\\
 u^{2+8\varepsilon}&\leq2^{h(u)}\leq2u^{2+8\varepsilon},
 \qquad 0<\varepsilon<1/256,\qquad T=\lfloor u^{399}\rfloor.
 \end{aligned}
 \tag{OPW.1}
\]
The rings $O_\ell$, their fraction fields $E_\ell$, ranks $b_\ell$,
full strict local conditions and derived indices are exactly those
of (FC.1)--(FC.2), (FL.2), (EU.2) and (TC.1).
In particular every coefficient prime remains in the original sum.

\paragraph{The new finite sets of actual triples.}
For each prime $\ell\geq5$ put
\[
 \begin{aligned}
 n_\ell&=v_2(\ell^2-1),\\
 c_\ell&=\min\{c\in\mathbb Z_{\geq0}:2^c>21\ell\,2^{2n_\ell}\},\\
 B(h;M)&=\max\{0,2^{h+1}-2^{M-2}\},
             \qquad C_\ell(h)=6B(h;c_\ell),\\
 \Pi(h)&=\{\ell\geq5:\ell\text{ prime},\ \ell\ne269,
                                      \ c_\ell\leq h+2\},\\
 P(h)&=\sum_{\ell\in\Pi(h)}C_\ell(h),\\
 C_3(h)&=B(h;8)=\max\{0,2^{h+1}-64\}.
 \end{aligned}
 \tag{OPW.2}
\]
Let $\Xi_7$ be all six characters modulo seven, with the trivial
character given its primitive conductor one.  Define
\[
 \begin{aligned}
 Z^{\geq5}(u)=\{\,&(\ell,(\delta\xi)|_{G_K},e):
       \ell\in\Pi(h(u)),\ \xi\in\Xi_7,\\
       &\delta\text{ primitive of conductor }2^c,
       \ c_\ell\leq c\leq h(u)+2,\ 1\leq e<T\},\\
 Z^{3,\mathrm{high}}(u)=\{\,&(3,\delta|_{G_K},e):
       \delta\text{ primitive of conductor }2^c,\\
       &8\leq c\leq h(u)+2,\ 1\leq e<T\},\\
 Z^{\mathrm{odd}}(u)&=Z^{\rm unif2}(u)
                 \sqcup Z^{\geq5}(u)\sqcup Z^{3,\mathrm{high}}(u).
 \end{aligned}
 \tag{OPW.3}
\]
Here $Z^{\rm unif2}$ is the entire previous union (TPW.3)--(TPW.5),
with all its old weights and exceptions unchanged.
The new labels have trivial cubic and auxiliary $q,d$ factors.
At coefficient three the conductor-seven factor is also trivial.
Both parities of the primitive two-power characters are included.
There is no new finite character at the coefficient prime.

For the old counts retain exactly the notation (TPW.2), (TPW.6):
\[
 \begin{aligned}
 N_T&=T-1,& B_T&=\lfloor(T-1)/268\rfloor,& A_T&=\lfloor T/67\rfloor,\\
 C(h)&=B(h;14)+B(h;17)+4B(h;19),\\
 R_T&=\sum_{r\in\{89,92,99,168,175,178\}}
                      \left\lfloor\frac{T+267-r}{268}\right\rfloor,\\
 S_T&=\sum_{\sigma=\pm1}\sum_{b=0}^1
       \sum_{r\in\mathcal F_{\sigma,b}}
                      \left\lfloor\frac{T+267-r}{268}\right\rfloor,
 \end{aligned}
 \tag{OPW.4}
\]
where $\mathcal F_{1,0}=\{243,248\}$,
$\mathcal F_{-1,0}=\{19,24\}$,
$\mathcal F_{1,1}=\{11,60,229\}$ and
$\mathcal F_{-1,1}=\{38,207,256\}$.
The symbols $B_T$ and $B(h;M)$ denote different functions.
Let $I_{399}^{\rm unifodd}$ omit precisely $Z^{\rm odd}$ from (TC.1).

\begin{theorem}[AC-G176-window: one union across growing coefficient primes]
\label{AC-G176-window}
Every triple in (OPW.3) is an actual original-coefficient strict zero.
The indicated unions are disjoint, and for positive integer $T$,
\[
 \begin{aligned}
 |Z^{\rm odd}\setminus Z^{\rm unif2}|&=(C_3(h)+P(h))N_T,\\
 |Z^{\rm odd}_{269}|&=(8+C(h))N_T-B_T+12A_T-R_T-S_T,\\
 |Z^{\rm odd}_{3}|&=(1+C_3(h))N_T,\\
 |Z^{\rm odd}_{\ell}|&=C_\ell(h)N_T\quad(\ell\in\Pi(h)),\\
 |Z^{\rm odd}|&=(9+C(h)+C_3(h)+P(h))N_T
                                     -B_T+12A_T-R_T-S_T.
 \end{aligned}
 \tag{OPW.5}
\]
No other coefficient prime has an omitted triple in this specified set.
The sum defining $P(h)$ is finite.  The sufficient condition
$84\ell(\ell+1)^2<2^{h+2}$ supplies a growing subrange of $\Pi(h)$,
apart from its explicit exclusion of $269$.
The inclusion of omitted sets is proper for $h\geq6$, $T\geq2$,
and hence eventually on every actual page.
Proper inclusion does not assert strict decrease of the numerical sum.

If $E_u$ is the normalized positive-index cost of the entire union, then
\[
 \begin{aligned}
 I_{399}&=I_{399}^{\rm unifodd}+E_u,\\
 I_{399}^{\rm unifodd}&\leq I_{399}^{\rm unif2}\leq I_{399},\\
 0\leq E_u&\leq8D_u^2T^2\log(8TD_u)
                  +\tfrac12D_u\mathcal R_uT^2+2D_ur_uc_0T\\
 &=O_\varepsilon\bigl(u^{802+16\varepsilon}(\log u)^2\bigr),
                         \qquad D_u=|G_u|.
 \end{aligned}
 \tag{OPW.6}
\]
These are nonnegative extended sums.  The error is used once for the
whole union, with no factor $|\Pi(h)|$, $P(h)$ or $C_3(h)$.
No individual positive derived index is asserted zero.
The new positive-relative-slack premise implies the original full
premise on the identical $W,J$, datum, page and prescribed pair,
after increasing only the initial threshold and halving slack.
Every older premise implies the new one by nonnegativity.
The actual source rank $R.t$, all simultaneous quantifiers and
the retained mixed realization $\mathcal R_{\rm cub}$ are unchanged.
\end{theorem}

\begin{proof}
\proofstep{1. Original coefficients at every newly used prime.}
The host contains $\mathbb Q(\zeta_{2^{h+2}},\zeta_7)$.
A primitive character of conductor $2^c$, $c\geq3$, has order
dividing $2^{c-2}$.  Products with characters modulo seven have
order dividing $3\cdot2^h$, the exact exponent in (FC.1).
Thus their values already belong to the prescribed $O_\ell$.
For each of the new labels the canonical scalar identity is
\[
 \begin{aligned}
 A_{\eta,e}
   &=(O_\ell\otimes_{\mathbb Z_\ell}\mathbb Q_\ell/\mathbb Z_\ell)
                  (\eta)(-e)\\
   &=(E_\ell/O_\ell)(\eta c_{\rm cyc}^{-e}).
 \end{aligned}
 \tag{OPW.7}
\]
The equality uses the same finite free $\mathbb Z_\ell$ module;
it does not extend its coefficient field or change its lattice.

For $\ell\in\Pi(h)$ apply (G175.3) with this $O_\ell$,
$s=0$, $\kappa=1$ and $\eta=\delta\xi$.
Its thresholds are exactly (OPW.2)--(OPW.3).
This includes $\ell=7$: the strict proof uses the least-common-multiple
class field, so the coincidence of the finite seven factor and the
coefficient prime is already accounted for.
For $\ell=3$ apply (G177.2) with the original $O_3$, $s=0$,
$\kappa=1$ and $\eta=\delta$.
That separate theorem works over $K_3=\mathbb Q_3(\mu_3)$ and
imposes no absolute-ramification restriction on $O_3$.
In particular it covers this actual ramified completion without
averaging a conductor-seven factor or a group of order divisible by three.
Both theorems assert the full strict zero at every positive weight,
including all residual weight classes.

Inflation from $G_{K,\Sigma_\ell}$ into $G_K$ is injective because
the original coefficient is unramified outside $\Sigma_\ell$.
An inflated (FC.2) cocycle kills inertia outside that set and
retains the same inside-set inertia conditions and full local zero
at every place over $\ell$.
It is therefore in the zero strict group of the appropriate theorem.
Equivalently, a zero inertia restriction for an unramified coefficient
is a zero cocycle; its conjugates and closed normal closure are
killed too.  No weakening to an unramified local condition is used.

\proofstep{2. Disjointness, finite prime range and exact counts.}
The old union has coefficient primes only three and $269$.
All triples in $Z^{\geq5}$ have a different coefficient prime.
The new coefficient-three labels are nontrivial, whereas the only
old coefficient-three label is trivial.  This proves disjointness
from the old union and between the two new families.
Distinct two-power and seven character factors remain distinct
on $G_K$: a rational character trivial there is either trivial or
$\chi_{-3}$, and a nontrivial character with conductor supported
only at two and seven cannot equal the latter.

At $269$ one has $n_{269}=3$ and $c_{269}=19$.
The AC-G175 labels there would all have $c\geq19$ and would
already belong to the AC-G173 part of $Z^{\rm unif2}$, whose
thresholds are $14,17,19$.  Excluding $269$ therefore loses
no newly supplied triple and avoids counting any old term twice.

There are $2^{c-2}$ primitive characters of conductor $2^c$.
For $M\geq3$ their sum over $M\leq c\leq h+2$ is
$B(h;M)$, including zero for the empty range.
There are six choices of $\xi$, so each prime in $\Pi(h)$
contributes $C_\ell(h)=6B(h;c_\ell)$ labels.
At coefficient three there is one choice of the trivial seven
factor, giving $C_3(h)=B(h;8)$.
Every label occurs at every one of the $N_T$ positive weights.
The old counts are (TPW.7), and adding the disjoint sets proves (OPW.5).

For every odd $\ell$, $n_\ell\geq3$.
If $c_\ell\leq h+2$, then
\[
       1344\ell\leq21\ell\,2^{2n_\ell}<2^{h+2}.
 \tag{OPW.8}
\]
This proves finiteness of $\Pi(h)$ without a prime-counting theorem.
The sufficient growing subrange follows from (G175.4).
For example $c_5=13$, so coefficient-five terms occur for $h\geq11$.
Already $C_3(h)>0$ for $h\geq6$; with $e=1$ this proves the
stated proper inclusion whenever $T\geq2$.

\proofstep{3. One local comparison over the complete enlarged union.}
Retain the original codifferent dual and rank $b_\ell$.
For each omitted triple, strict zero and (EU.2) give
\[
 \begin{aligned}
 \frac{\log\#k_{O_\ell}}{b_\ell}\max(0,i_{\ell,\eta,e})
 &\leq\frac{\log\#H^2(C_{\ell,\eta,e})}{b_\ell}\\
 &=\frac{|\log\#H^2(C_{\ell,\eta,e})
                  -\log\#\mathcal S_{\ell,\eta}(e)|}{b_\ell}.
 \end{aligned}
 \tag{OPW.9}
\]
At each fixed weight, the entire enlarged union is one subsum
of the nonnegative all-prime, all-character sum (FL.4).
That one bound already includes every eligible prime and every
eligible host character; neither growing number is a new multiplier.
Sum $16eD_u^2\log(8eD_u)+D_ue\mathcal R_u+2D_ur_uc_0$
once over $1\leq e<T$, using $\sum e\leq T^2/2$.
This gives (OPW.6) and the nonnegative additive decomposition,
without subtracting infinite quantities.
As before $D_u\leq288u^{2+8\varepsilon}$, $r_u\leq12$ and
$\mathcal R_u=O_\varepsilon(u^{\varepsilon/5}(\log u)^2)$.
With $T\leq u^{399}$ the result is
$O_\varepsilon(u^{802+16\varepsilon}(\log u)^2)$,
uniformly in the prescribed pair and little-$o$ of
$\exp(u^\varepsilon/2)$ for each fixed positive $\varepsilon$.

\proofstep{4. The identical window and the retained mixed input.}
Fix the retained finite datum, page and $\varepsilon$.
Assume $J\geq1$, $X_0$ and $s\in(0,1)$, independent of
$X,q,d,u,\ell,\eta,e$, give for every prescribed
$\operatorname{WinPair}(\varepsilon,X^{1/J},q,d)$, $X\geq X_0$,
one $W$ with $X^{1/J}\leq W$, $W^2\leq X$, and
\[
 I_{399}^{\rm unifodd}(H_u,u)\leq(1-s)\exp(u^\varepsilon/2)
                             \quad(W\leq u\leq W^2).
 \tag{OPW.10}
\]
Choose $U_0\geq2$ so (OPW.6) is at most
$(s/2)\exp(u^\varepsilon/2)$ for $u\geq U_0$.
Replace $X_0$ by $\max(X_0,U_0^J)$ and use exactly the already
supplied $W$.  The original full premise then holds on it with
slack $s/2$.
Uniform little-$o$ premises and fixed finite simultaneous families
are treated by the same estimate and a maximum of their uniform
error constants, without choosing separate windows.
Conversely, nonnegativity gives the new premise from every older one.
The arithmetic range still covers the actual $R.t$, and all field,
coefficient, character, weight and common-window quantifiers of
AC-K35 are retained.  Its mixed theta/Selmer and analytic realization
input $\mathcal R_{\rm cub}$ is not supplied or strengthened here.
\end{proof}

\paragraph{Remaining scope and geometric dictionary.}
All of the previous low-conductor exceptions remain.
At coefficient three the new theorem uses only pure two-power labels;
at primes at least five it allows the seven factor but no nontrivial
cubic or auxiliary $q,d$ factor.  Every unlisted triple, including
coefficient two and all unsupported prime/conductor combinations,
stays in the unknown sum.
No full numerical estimate, globally maximal classification,
mixed realization or H04.4 conclusion is asserted.
The rank-one finite characters, Tate tensor and original integral
cohomology use the existing AC-K25/26 dictionary, with the retained
SGA~1, Expos\'e~V, \S\S4,7 and Stacks~03PK scopes.
No new motive, archive attribution or Lean execution is claimed.
